\documentclass[10pt,a4paper]{amsart}
\usepackage{amsmath,amssymb,amsthm,graphicx}
\usepackage{verbatim}
\newenvironment{reference}{\comment}{\endcomment}
%\newenvironment{reference}{}{}
\newenvironment{slogan}{\comment}{\endcomment}
\newenvironment{history}{\comment}{\endcomment}

% For commutative diagrams we use Xy-pic
\usepackage[all]{xy}

% We use 2cell for 2-commutative diagrams.
\xyoption{2cell}
% Theorem environments.
%
\theoremstyle{plain}
\newtheorem{theorem}[subsection]{Theorem}
\newtheorem{proposition}[subsection]{Proposition}
\newtheorem{lemma}[subsection]{Lemma}

\theoremstyle{definition}
\newtheorem{definition}[subsection]{Definition}
\newtheorem{example}[subsection]{Example}
\newtheorem{exercise}[subsection]{Exercise}
\newtheorem{situation}[subsection]{Situation}

\theoremstyle{remark}
\newtheorem{remark}[subsection]{Remark}
\newtheorem{remarks}[subsection]{Remarks}

\numberwithin{equation}{subsection}

% Macros
%
\def\lim{\mathop{\mathrm{lim}}\nolimits}
\def\colim{\mathop{\mathrm{colim}}\nolimits}
\def\Spec{\mathop{\mathrm{Spec}}}
\def\Hom{\mathop{\mathrm{Hom}}\nolimits}
\def\Ext{\mathop{\mathrm{Ext}}\nolimits}
\def\SheafHom{\mathop{\mathcal{H}\!\mathit{om}}\nolimits}
\def\SheafExt{\mathop{\mathcal{E}\!\mathit{xt}}\nolimits}
\def\Sch{\mathit{Sch}}
\def\Mor{\mathop{\mathrm{Mor}}\nolimits}
\def\Ob{\mathop{\mathrm{Ob}}\nolimits}
\def\Sh{\mathop{\mathit{Sh}}\nolimits}
\def\NL{\mathop{N\!L}\nolimits}
\def\CH{\mathop{\mathrm{CH}}\nolimits}
\def\proetale{{pro\text{-}\acute{e}tale}}
\def\etale{{\acute{e}tale}}
\def\QCoh{\mathit{QCoh}}
\def\Ker{\mathop{\mathrm{Ker}}}
\def\Im{\mathop{\mathrm{Im}}}
\def\Coker{\mathop{\mathrm{Coker}}}
\def\Coim{\mathop{\mathrm{Coim}}}

% Boxtimes
%
\DeclareMathSymbol{\boxtimes}{\mathbin}{AMSa}{"02}

%
% Macros for moduli stacks/spaces
%
\def\QCohstack{\mathcal{QC}\!\mathit{oh}}
\def\Cohstack{\mathcal{C}\!\mathit{oh}}
\def\Spacesstack{\mathcal{S}\!\mathit{paces}}
\def\Quotfunctor{\mathrm{Quot}}
\def\Hilbfunctor{\mathrm{Hilb}}
\def\Curvesstack{\mathcal{C}\!\mathit{urves}}
\def\Polarizedstack{\mathcal{P}\!\mathit{olarized}}
\def\Complexesstack{\mathcal{C}\!\mathit{omplexes}}
% \Pic is the operator that assigns to X its picard group, usage \Pic(X)
% \Picardstack_{X/B} denotes the Picard stack of X over B
% \Picardfunctor_{X/B} denotes the Picard functor of X over B
\def\Pic{\mathop{\mathrm{Pic}}\nolimits}
\def\Picardstack{\mathcal{P}\!\mathit{ic}}
\def\Picardfunctor{\mathrm{Pic}}
\def\Deformationcategory{\mathcal{D}\!\mathit{ef}}
\usepackage{cleveref}
\crefname{originalequation}{original equation}{original equations}
\Crefname{originalequation}{Original equation}{Original equations}
\makeatletter
\newlabel{equation-list}{{048S}{source}{}{sourcepointer-1}{}}
\newlabel{equation-list@cref}{{[originalequation][0][]048S}{[1][0][]source}{}{}{}}
\makeatother
\crefname{originalreference}{original reference}{original references}
\Crefname{originalreference}{Original reference}{Original references}
\makeatletter
\newlabel{lemma-set-theoretic-pre-equivalence-geometric}{{0491}{source}{}{sourcepointer-2}{}}
\newlabel{lemma-set-theoretic-pre-equivalence-geometric@cref}{{[originalreference][0][]0491}{[1][0][]source}{}{}{}}
\makeatother
\crefname{originallemma}{original lemma}{original lemmas}
\Crefname{originallemma}{Original lemma}{Original lemmas}
\makeatletter
\newlabel{spaces-more-morphisms-lemma-when-universally-injective-radicial}{{0484}{source}{}{sourcepointer-3}{}}
\newlabel{spaces-more-morphisms-lemma-when-universally-injective-radicial@cref}{{[originallemma][0][]0484}{[1][0][]source}{}{}{}}
\makeatother
\crefname{originallemma}{original lemma}{original lemmas}
\Crefname{originallemma}{Original lemma}{Original lemmas}
\makeatletter
\newlabel{spaces-morphisms-lemma-universally-injective}{{040X}{source}{}{sourcepointer-4}{}}
\newlabel{spaces-morphisms-lemma-universally-injective@cref}{{[originallemma][0][]040X}{[1][0][]source}{}{}{}}
\makeatother
\title{Geometric point criterion for set-theoretic equivalence relations}
\author{The Stacks Project Authors}
\date{stacks-project:groupoids-quotients.tex:local-snapshot-20261012; GNU FDL 1.2 or later}
\begin{document}
\maketitle
\noindent\textit{Editorial note.} All delivered mathematical text is the author's own, copied from the pinned Stacks Project chapter file groupoids-quotients.tex; no statement or proof step was rewritten, reordered or removed.\par\medskip
\noindent\textit{Editorial note.} The counted proof is the author's complete argument: The refinement (4) requires one of decent, locally separated, or quasi-separated; without these the equivalence of (4) can fail.\par\medskip
\noindent\textit{Editorial note.} The label lemma-set-theoretic-equivalence-geometric names the counted statement, and the proof environment immediately following it in the source is delivered verbatim as its proof.\par\medskip
\noindent\textit{Editorial note.} The counted statement is carried by the author's own lemma environment in the chapter source.\par\medskip
\noindent\textit{Allowed standard prerequisites (editorial).}\begin{itemize}
\item Delivered card X2-alg-groupoids-quotients-set-theoretic-pre-equivalence-geometric (Lemma lemma-set-theoretic-pre-equivalence-geometric): the three-fold geometric criterion for j to be a set-theoretic pre-equivalence relation, including the containment of the images of the induced maps
\item Chapter Definition definition-set-theoretic-equivalence, included as a context span below: a set-theoretic equivalence relation is a universally injective set-theoretic pre-equivalence relation
\item Morphisms of Spaces, Lemma spaces-morphisms-lemma-universally-injective: universal injectivity can be tested by injectivity of the maps on k-valued points over algebraically closed fields
\item More on Morphisms of Spaces, Lemma spaces-more-morphisms-lemma-when-universally-injective-radicial: a universally injective morphism that is decent, or locally separated, or quasi-separated is radicial
\item Radicial morphisms: the fibres are spectra of purely inseparable field extensions, so over an algebraically closed field the reduction of a nonempty fibre is the spectrum of the base field itself
\end{itemize}\medskip
\ifdefined\hypertarget\hypertarget{sourcepointer-1}{}\fi
\noindent\textit{Source reference.} Original equation 048S: sources/stacks-groupoids-quotients/pinned-original-source.tex, source lines 813--813. Prerequisite pointer invoked at 1230; it names a standard fact of the same chapter that the counted argument uses, and every step of the counted proof built on it is retained verbatim in the delivered body.\par\medskip
\ifdefined\hypertarget\hypertarget{sourcepointer-2}{}\fi
\noindent\textit{Source reference.} Original reference 0491: sources/stacks-groupoids-quotients/pinned-original-source.tex, source lines 1091--1091. Prerequisite pointer invoked at 1246; it names a standard fact of the same chapter that the counted argument uses, and every step of the counted proof built on it is retained verbatim in the delivered body.\par\medskip
\ifdefined\hypertarget\hypertarget{sourcepointer-3}{}\fi
\noindent\textit{Source reference.} Original lemma 0484: sources/stacks-chow/tags.txt, source lines 5531--5531. Prerequisite pointer invoked at 1259; it names a standard fact of a sibling chapter that the counted argument uses, and every step of the counted proof built on it is retained verbatim in the delivered body.\par\medskip
\ifdefined\hypertarget\hypertarget{sourcepointer-4}{}\fi
\noindent\textit{Source reference.} Original lemma 040X: sources/stacks-chow/tags.txt, source lines 5272--5272. Prerequisite pointer invoked at 1250; it names a standard fact of a sibling chapter that the counted argument uses, and every step of the counted proof built on it is retained verbatim in the delivered body.\par\medskip
\begin{definition}
\label{definition-set-theoretic-equivalence}
Let $S$ be a scheme, and let $B$ be an algebraic space over $S$.
Let $j : R \to U \times_B U$ be a pre-relation over $B$.
\begin{enumerate}
\item We say $j$ is a {\it set-theoretic pre-equivalence relation} if
for all algebraically closed fields $k$ over $B$ the relation
$\sim_R$ on $U(k)$ defined by
$$
\overline{u} \sim_R \overline{u}'
\Leftrightarrow
\begin{matrix}
\exists\text{ field extension }K/k, \ \exists\ r \in R(K), \\
s(r) = \overline{u}, \ t(r) = \overline{u}'
\end{matrix}
$$
is an equivalence relation.
\item We say $j$ is a {\it set-theoretic equivalence relation}
if $j$ is universally injective and a set-theoretic pre-equivalence
relation.
\end{enumerate}
\end{definition}


\begin{lemma}
\label{lemma-set-theoretic-equivalence-geometric}
In the situation of Definition \ref{definition-set-theoretic-equivalence}.
The following are equivalent:
\begin{enumerate}
\item The morphism $j$ is a set-theoretic equivalence relation.
\item The morphism $j$ is universally injective and
$j(|R|) \subset |U \times_B U|$ contains the image of
$|j'|$ for any of the morphisms $j'$ as in Equation (\ref{equation-list}).
\item For every algebraically closed field $k$ over $B$ of sufficiently large
cardinality the map $j : R(k) \to U(k) \times U(k)$ is injective and
its image is an equivalence relation.
\end{enumerate}
If $j$ is decent, or locally separated, or quasi-separated
these are also equivalent to
\begin{enumerate}
\item[(4)] For every algebraically closed field $k$ over $B$
the map $j : R(k) \to U(k) \times U(k)$ is injective and its image
is an equivalence relation.
\end{enumerate}
\end{lemma}

\begin{proof}
The implications (1) $\Rightarrow$ (2) and (2) $\Rightarrow$ (3) follow from
Lemma \ref{lemma-set-theoretic-pre-equivalence-geometric}
and the definitions. The same lemma shows that (3) implies
$j$ is a set-theoretic pre-equivalence relation. But of course condition
(3) also implies that $j$ is universally injective, see
Morphisms of Spaces, Lemma \ref{spaces-morphisms-lemma-universally-injective},
so that $j$ is indeed a set-theoretic equivalence relation.
At this point we know that (1), (2), (3) are all equivalent.

\medskip\noindent
Condition (4) implies (3) without any further hypotheses on $j$. Assume $j$
is decent, or locally separated, or quasi-separated and the equivalent
conditions (1), (2), (3) hold. By
More on Morphisms of Spaces,
Lemma \ref{spaces-more-morphisms-lemma-when-universally-injective-radicial}
we see that $j$ is radicial.
Let $k$ be any algebraically closed field over $B$. Let
$\overline{u}, \overline{u}' \in U(k)$ with
$\overline{u} \sim_R \overline{u}'$. We see that
$R \times_{U \times_B U, (\overline{u}, \overline{u}')} \Spec(k)$
is nonempty. Hence, as $j$ is radicial, its reduction is the spectrum of a
field purely inseparable over $k$. As $k = \overline{k}$ we see that
it is the spectrum of $k$. Whence a point $\overline{r} \in R(k)$
with $t(\overline{r}) = \overline{u}$ and $s(\overline{r}) = \overline{u}'$
as desired.
\end{proof}

\end{document}
